\section{Justice Without Equivalence}

To call worship a debt of justice sounds, at first hearing, either too cold or too small. Cold, because justice evokes courts, ledgers, and settlement rather than awe before the living God. Small, because strict justice depends on equality, and between Creator and creature no equality exists. Yet St.\ Thomas Aquinas places the entire treatise on religion inside the treatise on justice in the \emph{Secunda Secundae} of the \emph{Summa Theologiae}. That placement rests on a deliberate realism: worship is not a decorative hobby for religious temperaments, nor an optional courtesy offered to a superior who might prefer something else, but an obligation grounded in the truth of what God is and what we are.

\subsection{Strict Justice and Its Potential Parts}

Aquinas begins with the classical Roman jurists' definition of justice: ``the perpetual and constant will to render to each one his right'' (\emph{ST} II-II, q.\ 58, a.\ 1).\endnote{St.\ Thomas Aquinas, \emph{Summa Theologiae} II-II, q.\ 58, a.\ 1, corpus (adapting Ulpian in \emph{Digest} 1.1.10); cf.\ II-II, q.\ 58, aa.\ 2 and 11.} Unlike temperance or fortitude, which order a person's own passions within himself, justice is always \emph{ad alterum}---directed toward another distinct person (\emph{ST} II-II, q.\ 58, a.\ 2). Its proper act is to establish equality of proportion in external deeds and things, rendering to another what is his own (\emph{ST} II-II, q.\ 58, a.\ 11). If I borrow a sum of money or purchase a house, commutative justice demands an equal return; once that equivalent is handed over, the debt is extinguished and the balance stands at zero.

Plainly, not every real obligation between persons fits that template. A child receives existence, nurture, and formation from his parents; a citizen receives protection and civic heritage from his country; every human being receives existence itself, every power of soul and body, and every moment of continuance from God. If justice existed only where equal repayment is possible, then the deepest debts of human life would fall outside justice altogether.

Aquinas resolves this difficulty through the category of \emph{potential parts} annexed to a principal virtue. A virtue is annexed to justice when it shares justice's orientation toward another person (\emph{ad alterum}) while falling short of justice's full perfection in one of two ways: either by falling short of \emph{equality}, or by falling short of strict \emph{legal due}. Three great virtues fall short of equality, and Aquinas names religion first among them---citing the very verse of the Psalter that asks what a creature can return to God:

\begin{quote}
For certain virtues there are which render another his due, but are unable to render the equal due. In the first place, whatever man renders to God is due, yet it cannot be equal, as though man rendered to God as much as he owes Him, according to Ps.\ 115:12, ``What shall I render to the Lord for all the things that He hath rendered to me?'' In this respect \emph{religion} is annexed to justice since, according to Tully (De invent.\ ii, 53), it consists in offering service and ceremonial rites or worship to ``some superior nature that men call divine.''\endnote{Aquinas, \emph{ST} II-II, q.\ 80, a.\ 1, corpus (English Dominican translation, Project Gutenberg eBook 18755).}
\end{quote}

Beneath religion stand \emph{piety} (\emph{pietas}), which renders service and deference to parents and country because no child can repay the gift of being begotten and reared, and \emph{observance} (\emph{observantia}), which honors those who excel in dignity or virtue. In all three, the impossibility of equivalence does not cancel the debt; it transforms the character of the debt. What cannot be paid in strict equality must still be answered according to the debtor's real capacity and the truth of the relationship. As Aquinas puts it in the next question, religion still observes a mean, not of ``absolute equality, because it is not possible to pay God as much as we owe Him, but equality in consideration of man's ability and God's acceptance'' (\emph{ST} II-II, q.\ 81, a.\ 5, ad 3). Nor does the creature's smallness diminish the dignity of the virtue: ``Virtue is praised because of the will, not because of the ability: and therefore if a man fall short of equality which is the mean of justice, through lack of ability, his virtue deserves no less praise, provided there be no failing on the part of his will'' (\emph{ST} II-II, q.\ 81, a.\ 6, ad 1).

This distinction saves worship from two opposite distortions. Against the distortion that reduces religion to arbitrary sentiment, annexation to justice establishes that God truly has a right to our acknowledgment; to withhold worship is not merely to miss an aesthetic experience, but to live in a falsehood toward the First Principle of our being. Against the distortion that turns religion into bargaining or self-justification, the impossibility of equivalence makes clear that no human offering ever places God in our debt or matches what we have received.

\subsection{Gratitude, Moral Debt, and the Unwilling Debtor}

The second group of virtues annexed to justice in \emph{ST} II-II, q.\ 80, a.\ 1 illuminates the same truth from another angle. These virtues fall short not primarily of equality, but of strict \emph{legal due} (\emph{debitum legale}), resting instead on \emph{moral due} (\emph{debitum morale})---that which a person owes out of the rectitude of virtue itself, even when no courtroom could enforce it. Among these stands \emph{gratitude} (\emph{gratia} or \emph{gratitudo}), which Aquinas treats directly in question 106.\endnote{Aquinas, \emph{ST} II-II, q.\ 106, aa.\ 1, 4, and 6.}

Aquinas orders the causes of debt in a descending scale: first and chiefly God, as the first principle of all our goods; second, parents, as proximate principles of life and upbringing; third, persons excelling in public dignity; fourth, particular human benefactors (\emph{ST} II-II, q.\ 106, a.\ 1). The higher always includes the lesser. Consequently, ``just as religion is superexcelling piety, so is it excelling thankfulness or gratitude'' (\emph{ST} II-II, q.\ 106, a.\ 1, ad 1).

Yet gratitude follows a logic completely foreign to commercial settlement. A legal debt must be discharged at once so that one does not keep another's property without consent. A debt of gratitude works differently, because a favor involves two things: the outward gift and the free benevolence of the giver's heart. Inwardly, gratitude responds immediately by receiving the favor graciously. Outwardly, however, a hurried attempt to match gift for gift at once destroys gratitude itself. Aquinas adopts Seneca's sharp diagnosis verbatim: ``he that wishes to repay too soon, is an unwilling debtor, and an unwilling debtor is ungrateful'' (\emph{ST} II-II, q.\ 106, a.\ 4).\endnote{Aquinas, \emph{ST} II-II, q.\ 106, a.\ 4, corpus (quoting Seneca, \emph{De Beneficiis} IV.40 and II.31).}

The person who rushes to repay a favor at the first possible moment is ungrateful because he cannot bear being indebted. He wants to convert a free gift into a completed transaction so that he may owe nothing more and regain his isolated autonomy. To accept a gift truly is to accept being bound to the giver in goodwill.

Aquinas pushes the insight even further when he asks whether gratitude should incline to return more than was received. Since a benefactor gave freely without compulsion, someone who returns merely an equivalent returns nothing freely of his own; he only gives back the principal. Gratitude therefore inclines to respond generously according to the disposition of the heart, even where equal return in kind is impossible, as between a child and his parents (\emph{ST} II-II, q.\ 106, a.\ 6, ad 1). And when an objector complains that this makes the debt of gratitude endless, Aquinas welcomes the consequence:

\begin{quote}
The debt of gratitude flows from charity, which the more it is paid the more it is due, according to Rom.\ 13:8, ``Owe no man anything, but to love one another.'' Wherefore it is not unreasonable if the obligation of gratitude has no limit.\endnote{Aquinas, \emph{ST} II-II, q.\ 106, a.\ 6, ad 2 (English Dominican translation, Project Gutenberg eBook 18755).}
\end{quote}

An unpayable debt feels like a prison only to someone who thinks debt is a humiliation to be erased. In the order of gift and love, an endless debt is not a crushing burden; it is the permanence of belonging, whereas the unwilling debtor wants only to pay God off and be left alone.

\subsection{Ingratitude as the Root of Idolatry}

What happens when a rational creature refuses this due return of acknowledgment and thanks? St.\ Paul traces the descent in the opening chapter of Romans, and the order of his verbs is telling:

\begin{quote}
Because that, when they knew God, they have not glorified him as God or given thanks: but became vain in their thoughts. And their foolish heart was darkened. For, professing themselves to be wise, they became fools. And they changed the glory of the incorruptible God into the likeness of the image of a corruptible man and of birds, and of fourfooted beasts and of creeping things. \dots\ Who changed the truth of God into a lie and worshipped and served the creature rather than the Creator, who is blessed for ever. (Rom 1:21--23, 25)
\end{quote}

What precedes idolatry is not simple ignorance of God, for the invisible attributes of the Creator---his eternal power and divinity---are intelligible through the things he has made (Rom 1:20). The rupture happens at the hinge of worship and thanksgiving: knowing God, they neither glorified him as God nor gave thanks (\emph{non sicut Deum glorificaverunt, aut gratias egerunt}). St.\ Augustine echoes this verse in \emph{City of God} X.3 when diagnosing the failure of the Platonists, who discerned the unchangeable Creator as the source of happiness yet withheld the sole worship due to him.\endnote{St.\ Augustine, \emph{De civitate Dei} X.3 (Marcus Dods translation, 1871; checked against Hoffmann, CSEL 40/1).}

The human person is not constructed to stand in neutral suspension, worshipping nothing. When gratitude toward the Creator is withheld, the religious impulse bends downward toward finite things---images, powers, status, appetite, or the self---and serves a lesser master that can neither create nor save.
